\chapter{Typhoid and paratyphoid fever}
\index{Typhoid and paratyphoid fever}

\section*{Definition and Overview}
Enteric fever is a severe, systemic infectious disease caused by specific serovars of the Gram-negative bacterium \textit{Salmonella enterica}. It encompasses both typhoid fever, caused by \textit{Salmonella enterica} serovar Typhi (\textit{S.} Typhi), and paratyphoid fever, caused by \textit{Salmonella enterica} serovars Paratyphi A, Paratyphi B, and Paratyphi C (\textit{S.} Paratyphi A, B, and C). While clinically indistinguishable, typhoid fever is generally more severe and prevalent than paratyphoid fever. Unlike non-typhoidal Salmonella species, which are zoonotic pathogens causing self-limiting localized gastroenteritis (food poisoning), the causative agents of enteric fever are human-restricted pathogens. This means humans are the sole reservoir of infection, and transmission occurs exclusively through the ingestion of food or water contaminated with the faeces or urine of infected individuals or chronic carriers.

Enteric fever is characterized by a prolonged, insidious onset of high-grade fever, headache, abdominal symptoms, and systemic toxaemia. The pathophysiological hallmark of the disease is the invasion of the gut-associated lymphoid tissue, followed by intracellular replication within the reticuloendothelial system (macrophages of the liver, spleen, bone marrow, and lymph nodes) and subsequent secondary bacteraemia. Historically, typhoid fever was a major cause of mortality in urban centers worldwide. A classic example of its transmission dynamics is the case of ``Typhoid Mary'' (Mary Mallon) in the early 20th century, an asymptomatic chronic carrier who worked as a cook and unwittingly infected dozens of people. The introduction of modern sanitation, municipal water chlorination, and sewage treatment systems in the late 19th and early 20th centuries led to a dramatic decline of the disease in high-income countries. However, in low- and middle-income countries (LMICs) where access to clean water and sanitation remains limited, enteric fever continues to pose a formidable public health challenge. Today, it also represents a significant concern in travel medicine, as travellers visiting endemic regions are at high risk of contracting the disease and importing it back to non-endemic countries. Furthermore, the rapid emergence and global spread of multidrug-resistant (MDR) and extensively drug-resistant (XDR) strains have compromised standard antibiotic therapies, rendering the management of enteric fever increasingly complex and costly.

\section*{Epidemiology}
The global epidemiology of enteric fever is characterized by a high burden in developing regions and a low, travel-associated incidence in developed countries. According to recent estimates by the World Health Organization (WHO) and global burden of disease studies, typhoid fever causes approximately 11 to 21 million illnesses and between 100,000 and 200,000 deaths annually worldwide. Paratyphoid fever accounts for an estimated additional 5 to 6 million cases per year. The vast majority of these infections occur in resource-limited settings characterized by inadequate sanitation, poor personal hygiene, and lack of access to safe drinking water.

Geographically, South Asia (including India, Pakistan, Bangladesh, and Nepal) is the highest-burden region globally, with incidence rates often exceeding 100 cases per 100,000 population per year. Other highly endemic areas include Southeast Asia, Sub-Saharan Africa (particularly parts of East and West Africa), and pockets of Central and South America. In these endemic regions, enteric fever is primarily a disease of children and young adults, with the peak incidence occurring in school-aged children between 2 and 15 years of age. Infants and toddlers under the age of 2 can also be affected, often presenting with atypical or severe clinical manifestations. The high burden in young populations is attributed to a combination of immature immune responses, lack of prior exposure, and behaviors that increase exposure to contaminated food and water (such as eating from street vendors).

In contrast, in high-income regions such as North America, Western Europe, and some countries, enteric fever is rare, with an annual incidence of less than 15 cases per 100,000 population. Most cases diagnosed in these areas are ``imported'' by international travellers returning from endemic regions, or occur among immigrants visiting friends and relatives (VFR travellers), who are often less likely to seek pre-travel vaccination. A small proportion of cases in developed nations can be traced to localized outbreaks linked to asymptomatic chronic carriers working in the food industry.

Epidemiological patterns are also influenced by seasonal variations. In many endemic regions, the incidence of enteric fever peaks during the hot, dry season (when water scarcity leads to the use of unsafe water sources) or immediately following the onset of the monsoon rains (when flooding causes sewage to contaminate municipal and well water supplies). Understanding these dynamics is crucial for planning targeted public health interventions, such as vaccination campaigns and infrastructure development.

\section*{Aetiology and Pathophysiology}
Enteric fever is caused by host-restricted, Gram-negative, flagellated, non-spore-forming, facultative anaerobic bacilli belonging to the family Enterobacteriaceae. The primary pathogen is \textit{Salmonella enterica} subspecies \textit{enterica} serovar Typhi. Paratyphoid fever is caused by serovars Paratyphi A, Paratyphi B (also known as \textit{S. enterica} serovar Schottmuelleri), and Paratyphi C (also known as \textit{S. enterica} serovar Hirschfeldii). These organisms are highly adapted to the human host and possess a complex array of virulence factors that allow them to evade host immune mechanisms and establish systemic infection.

The clinical and pathological progression of enteric fever is defined by a series of distinct pathophysiological stages:
\begin{itemize}
\item \textbf{Ingestion and Gastric Transit:} Transmission occurs via the faecal-oral route through the ingestion of food or water contaminated with human faeces or urine. The infectious dose of \textit{S.} Typhi is relatively low, ranging from $10^3$ to $10^5$ organisms. Upon ingestion, the bacteria must survive the acidic environment of the stomach. \textit{S.} Typhi possesses an acid tolerance response (ATR) that enables it to adjust its internal pH and survive gastric acidity, especially when ingested with food or large volumes of liquid that buffer stomach acid.
\item \textbf{Mucosal Invasion and Intestinal Penetration:} Once they reach the distal ileum, the bacteria attach to and invade the specialized microfold (M) cells of the follicle-associated epithelium overlying the Peyer's patches. They can also directly invade enterocytes. This invasion is mediated by a type III secretion system (T3SS-1) encoded by Salmonella Pathogenicity Island 1 (SPI-1), which injects effector proteins into the host cells, triggering actin cytoskeletal rearrangements, membrane ruffling, and subsequent internalization of the bacteria.
\item \textbf{Intracellular Survival within Macrophages:} After crossing the epithelial barrier, the bacteria are phagocytosed by macrophages, dendritic cells, and polymorphonuclear leukocytes in the lamina propria. Instead of being destroyed, \textit{S.} Typhi survives and replicates within these cells. It utilizes a second type III secretion system (T3SS-2), encoded by Salmonella Pathogenicity Island 2 (SPI-2), to prevent the fusion of the phagosome with lysosomes, creating a specialized intracellular niche known as the Salmonella-containing vacuole (SCV). The Vi (virulence) capsular polysaccharide, which is unique to \textit{S.} Typhi and some strains of \textit{S.} Paratyphi C, plays a critical role in evading host immune recognition by masking lipopolysaccharide (LPS) O-antigens, inhibiting neutrophil chemotaxis, and preventing complement-mediated lysis.
\item \textbf{Lymphatic Dissemination and Primary Bacteraemia:} The infected macrophages migrate via the lymphatic system to the mesenteric lymph nodes. From there, the bacteria enter the systemic circulation via the thoracic duct, resulting in a transient, primary, low-grade bacteraemia. During this early phase (which corresponds to the incubation period), the patient is typically asymptomatic or has mild, non-specific symptoms.
\item \textbf{Reticuloendothelial Sequestration and Multiplication:} The circulating bacteria are rapidly cleared from the bloodstream by the mononuclear phagocyte system and sequestered within the reticuloendothelial organs, including the liver, spleen, bone marrow, and lymph nodes. Within these organs, the bacteria undergo intensive intracellular multiplication over a period of 7 to 14 days.
\item \textbf{Secondary Bacteraemia and Clinical Disease:} Once a critical bacterial density is reached, the pathogens are released from the reticuloendothelial tissues back into the bloodstream in massive numbers. This secondary, sustained bacteraemia marks the onset of clinical illness. The release of bacterial endotoxins (LPS) and the host's release of pro-inflammatory cytokines (such as tumor necrosis factor-alpha, interleukin-1, and interleukin-6) in response to the bacteraemia trigger the characteristic clinical symptoms, including sustained high fever, headache, malaise, and abdominal pain.
\item \textbf{Gallbladder Colonization and Reinvasion of the Gut:} During the secondary bacteraemia, the bacteria travel to the gallbladder via the bloodstream or the biliary tract. \textit{S.} Typhi is highly resistant to the bactericidal effects of bile and multiplies rapidly within the gallbladder. The bacteria are subsequently excreted in the bile back into the intestinal lumen. This leads to the reinvasion of the Peyer's patches in the distal ileum, which are now sensitized from the primary infection. This secondary invasion triggers an intense inflammatory reaction, characterized by mononuclear cell infiltration, hyperplasia, necrosis, and ulceration of Peyer's patches, which can lead to severe intestinal complications.
\item \textbf{Development of the Chronic Carrier State:} In approximately 2\% to 5\% of infected individuals, the bacteria fail to be cleared and establish a chronic, persistent infection in the gallbladder. \textit{S.} Typhi can form biofilms on the surface of gallstones, which protects the bacteria from both host immune defenses and antibiotic therapy. These chronic carriers continuously or intermittently shed the bacteria in their faeces for more than a year, serving as the primary epidemiologic reservoir for the disease.
\end{itemize}

\section*{Clinical Features}
The clinical presentation of enteric fever is highly variable, ranging from a mild illness with low-grade fever to a severe, life-threatening multi-organ disease. The incubation period typically lasts 7 to 14 days, but can range from 3 to 60 days depending on the infectious dose, host immune status, and prior vaccination. Classically, in the absence of effective antibiotic therapy, the disease progresses through four distinct stages, each lasting approximately one week:
\begin{itemize}
\item \textbf{First Week (Prodromal Stage):} The onset of illness is characteristically insidious. The primary symptom is a progressive, step-ladder rise in body temperature, which increases daily to reach a sustained plateau of 39$^\circ$C to 40$^\circ$C. This is accompanied by a dull, generalized headache, dry cough, malaise, sore throat, and muscle aches (myalgia). Abdominal discomfort, anorexia, and constipation (especially in adults) or mild diarrhoea (more common in children) are frequent early features. A classic clinical sign is relative bradycardia (Faget's sign), in which the heart rate is paradoxically slow relative to the height of the fever (e.g., a pulse of 80 beats per minute during a fever of 40$^\circ$C).
\item \textbf{Second Week (Fastigium Stage):} During this week, the fever remains high and plateaued. The patient appears toxic and severely exhausted. Neuropsychiatric symptoms, collectively referred to as the ``typhoid state,'' may begin to manifest; these include extreme apathy, confusion, delirium, or stupor. Abdominal symptoms become more pronounced, with significant abdominal distension, diffuse tenderness, and hepatosplenomegaly. Around the 7th to 10th day, ``rose spots'' appear in 5\% to 30\% of patients. These are faint, salmon-coloured, slightly raised macules, 2 to 4 millimetres in diameter, which blanch under pressure. They typically appear in crops of 5 to 15 lesions on the abdomen and lower chest, lasting for 2 to 3 days before resolving without scarring. Diarrhoea described as having a ``pea-soup'' consistency and appearance may occur, accompanied by significant abdominal gurgling and flatulence.
\item \textbf{Third Week (Stage of Complications):} This is the most critical phase of the untreated illness, during which the patient's nutritional and physical state deteriorates severely. The patient enters a state of ``coma vigil,'' lying motionless with eyes half-open, showing profound apathy and muttering delirium. The high fever persists, and signs of severe toxaemia are prominent. The risk of life-threatening intestinal complications peaks during this week due to the necrosis and ulceration of Peyer's patches. The patient may develop sudden, massive gastrointestinal haemorrhage or intestinal perforation, presenting as acute abdomen and septic shock. Other complications, such as myocarditis, toxic encephalopathy, and secondary bacterial infections, may also emerge.
\item \textbf{Fourth Week (Defervescence and Convalescence):} If the patient survives the third week, a slow recovery begins. The fever gradually declines over several days (defervescence) in a step-ladder fashion, and the appetite and mental status begin to normalize. However, the patient remains profoundly weak and emaciated, and complete recovery can take several weeks or months. Relapse of symptoms can occur in 5\% to 10\% of patients during the convalescent phase, typically 1 to 3 weeks after the resolution of the initial fever.
\end{itemize}
It is important to note that the classic, four-week course is rarely seen today in clinical practice due to the prompt administration of effective antibiotic therapy, which usually aborts the progression of the disease within 3 to 5 days. Additionally, atypical presentations are common, particularly in young children, who may present with simple gastroenteritis, isolated respiratory symptoms, or acute convulsive seizures.

\section*{Differential Diagnosis}
Because the early clinical features of enteric fever are non-specific and overlap significantly with other acute febrile illnesses, establishing a differential diagnosis is essential, especially in endemic areas and returning travellers. The primary conditions that must be considered include:
\begin{itemize}
\item \textbf{Malaria:} Malaria is the most important differential diagnosis for any febrile patient in or returning from the tropics. Unlike the insidious onset of typhoid, malaria often presents with sudden onset of high fever, cyclic paroxysms (chills, spiking fever, and profuse sweating), severe rigors, and anaemia. While splenomegaly is common to both, the diagnosis of malaria is rapidly confirmed by thick and thin blood films or rapid diagnostic tests (RDTs) for Plasmodium antigens.
\item \textbf{Dengue Fever:} Dengue is characterized by a sudden onset of high fever, severe retro-orbital headache, retro-orbital pain, and intense muscle and joint pain (``breakbone fever''). In contrast to the faint rose spots of typhoid, dengue presents with a flushing facial rash early on, followed by a maculopapular, pruritic rash later in the course. Significant leukopenia and thrombocytopenia are common. The absence of a step-ladder prodrome and the presence of severe arthralgia help differentiate it from typhoid.
\item \textbf{Rickettsial Infections (Scrub Typhus and Murine Typhus):} Rickettsial diseases present with high fever, headache, myalgia, and a macular or maculopapular rash. A key diagnostic feature is the presence of an eschar (a painless ulcer with a black necrotic crust) at the site of the vector bite, along with localized or generalized lymphadenopathy. Diagnostic serology or PCR for rickettsiae can differentiate these from enteric fever.
\item \textbf{Leptospirosis:} Leptospirosis should be suspected in patients with a history of exposure to water or soil contaminated with animal urine (e.g., during farming, flooding, or water sports). It presents with sudden high fever, severe headache, intense muscle tenderness (particularly in the calves and lumbar region), conjunctival suffusion (redness of the eyes without discharge), and in severe cases (Weil's disease), jaundice and renal failure.
\item \textbf{Brucellosis:} Brucellosis is a zoonotic infection contracted through contact with infected animals or ingestion of unpasteurized dairy. It causes a chronic, undulant (wavelike) fever, profound drenching sweats with a characteristically sour or moldy odour, arthralgia, and body aches. It lacks the acute intestinal complications (such as perforation) seen in typhoid and has a more prolonged, relapsing course.
\item \textbf{Amebic Liver Abscess:} This condition presents with high fever and right upper quadrant abdominal pain. Physical examination reveals tender hepatomegaly, and laboratory tests show marked neutrophilic leukocytosis. Unlike typhoid, blood cultures are negative for \textit{S.} Typhi, and abdominal ultrasound or CT scan reveals one or more space-occupying fluid collections in the liver.
\item \textbf{Acute Viral Hepatitis (Hepatitis A and E):} These enterically transmitted viral infections present with prodromal fever, nausea, vomiting, and anorexia, followed by marked jaundice, dark urine, and pale stools. While typhoid can cause mild hepatitis, viral hepatitis is distinguished by extremely elevated serum transaminases (ALT and AST often exceeding 1000 U/L) and positive serology for IgM anti-HAV or IgM anti-HEV.
\item \textbf{Inflammatory Bowel Disease (IBD):} Crohn's disease or ulcerative colitis can present with chronic fever, abdominal pain, and diarrhoea. However, they lack the acute bacteraemic onset, relative bradycardia, and rose spots of typhoid, and are characterized by a chronic relapsing history, specific endoscopic findings, and negative cultures.
\end{itemize}

\section*{When to Seek Medical Care}
Enteric fever is a potentially life-threatening infection that requires prompt medical evaluation and treatment. Delaying care increases the risk of severe complications, such as intestinal perforation, haemorrhage, and septic shock.
Medical attention should be sought immediately in the following circumstances:
\begin{itemize}
\item \textbf{Travel-Related Fever:} Anyone who develops a persistent fever (38$^\circ$C or higher) during or within three weeks of returning from an area where typhoid is common (such as South Asia, Southeast Asia, or Sub-Saharan Africa) must be evaluated by a healthcare professional immediately.
\item \textbf{Progressive Febrile Illness:} A fever that starts low and steadily climbs over several days, particularly if accompanied by a severe headache, extreme fatigue, dry cough, and abdominal pain, warrants prompt medical assessment.
\item \textbf{Exposure to a Known Case:} Individuals who have been in close contact with someone diagnosed with typhoid or paratyphoid fever and who develop any symptoms of infection should seek medical testing.
\end{itemize}

\noindent Certain clinical signs indicate severe, life-threatening complications. Emergency medical services must be contacted immediately (e.g., by calling \textbf{local emergency services} in many countries, \textbf{999} in the United Kingdom, \textbf{911} in the United States and Canada, or \textbf{112} in the European Union) if a patient exhibits any of the following ``red flag'' symptoms:
\begin{itemize}
\item Sudden, severe, and worsening abdominal pain, which may be accompanied by rigidity, extreme tenderness, or abdominal swelling (signs of intestinal perforation and peritonitis).
\item Passing blood in the stool, black tarry stools (melaena), or vomiting blood or coffee-ground-like material (signs of severe gastrointestinal haemorrhage).
\item Altered mental status, including confusion, delirium, extreme lethargy, stupor, or loss of consciousness (signs of typhoid encephalopathy or septic shock).
\item Difficulty breathing, rapid breathing, or chest pain.
\item Signs of shock, such as cold, clammy skin, bluish lips or fingernails, rapid and weak pulse, dizziness, or fainting.
\item Persistent vomiting and inability to keep oral fluids or medications down, leading to severe dehydration.
\end{itemize}

\section*{Diagnosis and Investigations}
The diagnosis of enteric fever requires a combination of clinical suspicion, detailed travel and exposure history, and confirmatory laboratory investigations. Clinical diagnosis alone is unreliable due to the non-specific nature of the symptoms.
The key diagnostic investigations include:
\begin{itemize}
\item \textbf{Microbiological Cultures (Gold Standard):} The isolation of \textit{S.} Typhi or \textit{S.} Paratyphi from clinical specimens remains the definitive method of diagnosis.
  \begin{itemize}
  \item \textit{Blood Culture:} Blood cultures are the primary diagnostic tool and are positive in 60\% to 80\% of patients during the first week of illness. The sensitivity declines to 30\% to 40\% by the third week as the bacteraemia subsides. The volume of blood collected is critical; at least 10 to 15 millilitres is required for adults, and 2 to 4 millilitres for children, to optimize yield.
  \item \textit{Bone Marrow Aspirate Culture:} This is the most sensitive diagnostic method, with a positive yield of 85\% to 95\%. Its sensitivity is not significantly reduced by prior antibiotic therapy (unlike blood cultures). Because it is an invasive and painful procedure, it is typically reserved for patients with suspected enteric fever who have negative blood cultures and do not respond to empirical therapy.
  \item \textit{Stool and Urine Cultures:} Stool cultures are positive in approximately 30\% to 50\% of patients during the second and third weeks of illness. They are valuable for diagnosing patients who present late, monitoring for clearance after treatment, and screening food handlers or identifying chronic carriers. Urine cultures are positive in about 10\% of cases, typically in the third week.
  \item \textit{Duodenal String Test:} This involves swallowing a gelatin capsule containing a string to sample duodenal secretions and bile. It has a high sensitivity for culturing \textit{S.} Typhi but is rarely used in routine clinical practice due to patient discomfort.
  \end{itemize}
\item \textbf{Serological Testing:}
  \begin{itemize}
  \item \textit{Widal Test:} The Widal test measures agglutinating antibodies against the lipopolysaccharide somatic O antigen and flagellar H antigen of \textit{S.} Typhi. It is historically widely used but has significant limitations, including high rates of false-positive results (due to cross-reactivity with other Salmonella serovars or non-specific immune stimulation) and false-negative results (especially if performed too early in the disease). A single test is rarely diagnostic; a four-fold rise in antibody titres between acute and convalescent sera collected 10 to 14 days apart is required for confirmation, which limits its utility for acute clinical management.
  \item \textit{Rapid Diagnostic Tests (RDTs):} Commercial assays (such as Typhidot and Tubex) detect IgM and IgG antibodies directed against specific outer membrane proteins of \textit{S.} Typhi. IgM detection indicates acute infection. While they provide rapid results (within hours), their sensitivity and specificity are variable, and they cannot replace cultures for definitive diagnosis.
  \end{itemize}
\item \textbf{Molecular Testing (PCR):} Polymerase Chain Reaction assays designed to amplify specific genes of \textit{S.} Typhi (such as the \textit{fliC} flagellin gene or \textit{viaB} capsular gene) offer high sensitivity and specificity. PCR can detect bacterial DNA in blood or stool within hours, even in patients who have received antibiotics. However, high cost and the need for specialized laboratory equipment limit its availability in resource-constrained endemic areas.
\item \textbf{General Supportive Investigations:}
  \begin{itemize}
  \item \textit{Complete Blood Count (CBC):} Common findings include mild anaemia, thrombocytopenia, and leukopenia (white blood cell count typically between 3,000 and 5,000 per microlitre) with relative lymphocytosis. This contrasts with most other bacterial infections, which cause leukocytosis. A marked leukocytosis in a typhoid patient suggests a complication such as intestinal perforation or secondary localized infection.
  \item \textit{Liver Function Tests (LFTs):} Mild to moderate elevations in serum transaminases (alanine aminotransferase [ALT] and aspartate aminotransferase [AST]) and bilirubin are common, reflecting reactive hepatitis.
  \item \textit{Inflammatory Markers:} C-reactive protein (CRP) and erythrocyte sedimentation rate (ESR) are routinely elevated, though they are non-specific.
  \end{itemize}
\end{itemize}

\section*{Management}
The management of enteric fever involves a combination of definitive antimicrobial therapy, supportive care, monitoring for complications, and in select cases, surgical intervention. The primary goal of treatment is to eradicate the infection, alleviate symptoms, and prevent life-threatening complications.

\begin{itemize}
\item \textbf{Antimicrobial Therapy:} Antibiotics are the cornerstone of treatment and must be initiated promptly upon clinical suspicion. The choice of antibiotic depends on the severity of the illness, the geographic origin of the infection, and local antibiotic resistance patterns:
  \begin{itemize}
  \item \textit{The Challenge of Drug Resistance:} Historically, chloramphenicol, ampicillin, and co-trimoxazole were the standard first-line therapies. However, the emergence of Multidrug-Resistant (MDR) typhoid (resistant to all three classes) in the late 1980s led to a shift toward fluoroquinolones (such as ciprofloxacin). In recent decades, fluoroquinolone-resistant strains have become widespread, particularly in South Asia. Furthermore, Extensively Drug-Resistant (XDR) \textit{S.} Typhi (which is resistant to first-line drugs, fluoroquinolones, and third-generation cephalosporins) has emerged in Pakistan and other regions, leaving very limited oral treatment options.
  \item \textit{Empirical Treatment for Uncomplicated Typhoid:} In areas with high fluoroquinolone resistance, oral azithromycin (20 mg/kg/day, up to 1 gram daily, for 7 days) is the preferred treatment for uncomplicated cases, showing excellent efficacy and low relapse rates. Oral cefixime (15-20 mg/kg/day, for 10-14 days) is an alternative.
  \item \textit{Treatment for Severe or Complicated Typhoid:} For patients with severe illness, persistent vomiting, or suspected complications, intravenous therapy is required. Intravenous ceftriaxone (2 grams daily for adults; 80 mg/kg/day for children, for 10 to 14 days) is the standard empirical choice. In cases of confirmed or suspected XDR typhoid, intravenous carbapenems (such as meropenem, 1-2 grams every 8 hours) or tigecycline are used.
  \item \textit{Susceptible Strains:} If the strain is confirmed to be fully susceptible, ciprofloxacin (500 mg orally twice daily or 400 mg IV every 12 hours for 7 to 10 days) remains an efficient and cost-effective option.
  \end{itemize}
\item \textbf{Adjuvant Corticosteroids:} High-dose systemic corticosteroids are recommended for patients presenting with severe, complicated enteric fever characterized by neuropsychiatric manifestations (delirium, stupor, coma) or septic shock. Intravenous dexamethasone (initial dose of 3 mg/kg, followed by 1 mg/kg every 6 hours for a total of 8 doses) has been shown to reduce mortality significantly in these severe cases. It should only be given in conjunction with effective antibiotic coverage.
\item \textbf{Supportive Care and Symptom Management:}
  \begin{itemize}
  \item \textit{Fluid and Electrolyte Replacement:} Maintaining hydration is vital. Oral rehydration salts (ORS) should be given to patients with mild-to-moderate dehydration. Intravenous fluids (such as normal saline or Ringer's lactate) are required for patients with severe dehydration, persistent vomiting, or shock.
  \item \textit{Antipyretics:} Paracetamol (acetaminophen) should be used to control high fever and headache. Aspirin and nonsteroidal anti-inflammatory drugs (NSAIDs like ibuprofen) must be strictly avoided, as they can irritate the gastric mucosa and increase the risk of gastrointestinal bleeding and platelet dysfunction.
  \item \textit{Nutritional Support:} Patients should be encouraged to consume a soft, easily digestible, high-calorie diet as tolerated. Strict fasting is unnecessary and can worsen muscle wasting and delay recovery.
  \end{itemize}
\item \textbf{Surgical Management:} Intestinal perforation is a surgical emergency. Immediate exploratory laparotomy is required, accompanied by broad-spectrum antibiotic coverage for polymicrobial peritonitis (e.g., adding metronidazole or gentamicin to ceftriaxone). The surgical procedure typically involves primary closure of the perforation in one or two layers, or segment resection of the ileum if multiple perforations or extensive necrosis are present. Thorough peritoneal lavage is critical.
\end{itemize}

\section*{Complications}
In the absence of prompt antibiotic treatment, approximately 10\% to 20\% of patients with enteric fever develop serious, potentially life-threatening complications. These typically emerge during the third week of illness and are related to the intense inflammatory response and tissue necrosis at the sites of bacterial replication.
The major complications include:
\begin{itemize}
\item \textbf{Intestinal Haemorrhage:} Occurs in 2\% to 10\% of cases, resulting from the erosion of blood vessels as Peyer's patches ulcerate. It can manifest as minor, occult bleeding detected in the stool, or as sudden, massive haematochezia (bright red blood in stool) or melaena (black, tarry stools). Severe haemorrhage presents with a sudden drop in body temperature and blood pressure, rapid heart rate, and signs of hypovolaemic shock, requiring urgent blood transfusion and supportive care.
\item \textbf{Intestinal Perforation:} This is the most dreaded complication, occurring in 1\% to 3\% of hospitalized patients, typically in the terminal ileum. Perforation allows intestinal contents to spill into the peritoneal cavity, causing acute chemical and bacterial peritonitis. It presents clinically as a sudden, sharp, localized pain in the right lower quadrant of the abdomen, followed by generalized abdominal rigidity, guarding, rebound tenderness, and absent bowel sounds. The patient's condition rapidly deteriorates with high fever, vomiting, rapid shallow breathing, hypotension, and septic shock. Plain abdominal X-rays showing free air under the diaphragm (pneumoperitoneum) support the diagnosis.
\item \textbf{Typhoid Encephalopathy and Neuropsychiatric Complications:} Severe toxaemia can lead to central nervous system involvement. This presents as a ``typhoid state'' (coma vigil), muttering delirium, confusion, hallucinations, stupor, or coma. Other neurological complications include meningitis, cerebellar ataxia, Guillain-Barr\'e syndrome, and peripheral neuritis.
\item \textbf{Cardiovascular Complications:} Myocarditis is a significant complication caused by bacterial endotoxins. It can manifest as electrocardiographic abnormalities (such as prolonged PR interval, T-wave inversion, or arrhythmias), cardiomegaly, or acute congestive heart failure. Toxic shock and peripheral vascular collapse can also occur.
\item \textbf{Hepatobiliary Complications:} Acute cholecystitis can develop, presenting with right upper quadrant pain, tenderness, and vomiting. Typhoid hepatitis, characterized by marked jaundice, hepatomegaly, and significantly elevated liver enzymes, can mimic viral hepatitis.
\item \textbf{Focal Infections and Abscesses:} Bacteria can seed into various organs via the bloodstream, leading to localized infections. Examples include osteomyelitis (particularly in patients with sickle cell disease or other haemoglobinopathies), septic arthritis, splenic abscess, renal abscess, and pneumonia.
\item \textbf{Chronic Carrier State:} About 2\% to 5\% of treated or untreated patients fail to clear \textit{S.} Typhi from their systems and continue to excrete the bacteria in their faeces or urine for more than 12 months. This chronic carriage is more common in women, older adults, and individuals with pre-existing biliary tract pathology (such as gallstones or chronic cholecystitis). Chronic carriers are epidemiologically significant as reservoirs for transmission and have a significantly increased risk of developing gallbladder cancer.
\end{itemize}

\section*{Prognosis and Outcomes}
The prognosis of enteric fever has been dramatically transformed by the introduction of antimicrobial therapy. In the pre-antibiotic era, typhoid fever was a highly lethal disease, with mortality rates ranging from 10\% to 20\%, and even higher in vulnerable populations such as infants and the elderly. Today, with prompt diagnosis and administration of appropriate antibiotics, the prognosis is excellent:
\begin{itemize}
\item \textbf{Mortality Rate:} In patients receiving timely and effective antibiotic treatment, the mortality rate is less than 1\%. However, in resource-poor settings where access to healthcare is delayed, or in areas with high rates of XDR \textit{S.} Typhi where empirical treatment fails, mortality rates can still exceed 5\% to 10\%.
\item \textbf{Recovery Time:} Once effective antibiotic therapy is commenced, the patient's fever typically resolves within 3 to 5 days, and systemic symptoms begin to improve. Complete clinical recovery and return to baseline strength and energy levels can take 2 to 4 weeks. Without antibiotics, recovery is slow, taking 4 to 8 weeks, during which severe muscle wasting and weakness persist.
\item \textbf{Relapse Rates:} Relapse occurs in approximately 5\% to 10\% of patients, typically 1 to 3 weeks after the discontinuation of antibiotic therapy. The clinical features of a relapse are identical to the primary infection but are usually milder and of shorter duration. Relapse is not due to acquired drug resistance; the relapsing bacterial strain is almost always susceptible to the same antibiotic used initially, and a second course of treatment is effective.
\item \textbf{Factors Affecting Prognosis:} Poor prognostic factors include delayed initiation of effective antibiotics, infection with drug-resistant strains (MDR or XDR), presence of severe complications (such as intestinal perforation, haemorrhage, or encephalopathy), extremes of age (infants under 1 year and elderly patients), pre-existing chronic medical conditions, and severe malnutrition.
\end{itemize}

\section*{Prevention and Self-Care}
Prevention of enteric fever is achieved through a multi-faceted approach combining public health measures, safe food and water practices, vaccination, and proper self-care during and after infection. Because the disease is transmitted via the faecal-oral route, interrupting transmission pathways is key to control.

\begin{itemize}
\item \textbf{Water, Sanitation, and Hygiene (WASH):} This is the most fundamental and sustainable method of preventing enteric fever:
  \begin{itemize}
  \item Access to clean, municipally treated, chlorinated drinking water.
  \item Proper disposal and treatment of human sewage to prevent contamination of water systems.
  \item Frequent handwashing with soap and clean, running water, especially after using the toilet, changing nappies, and before preparing or consuming food.
  \end{itemize}
\item \textbf{Safe Food Practices for Travellers:} When travelling to endemic areas, individuals should adhere to strict food safety guidelines, summarized by the phrase ``Cook it, peel it, boil it, or forget it'':
  \begin{itemize}
  \item Drink only bottled, boiled, or chemically treated water. Avoid ice unless it is made from safe water.
  \item Avoid raw fruits and vegetables unless they can be peeled by the consumer.
  \item Consume only hot, freshly cooked food. Avoid street food and raw or undercooked seafood.
  \item Avoid unpasteurized milk and dairy products.
  \end{itemize}
\item \textbf{Vaccination:} Immunization is a highly effective preventive measure, recommended for travellers to endemic regions, laboratory personnel working with the bacteria, and as part of national immunization programs in high-burden countries:
  \begin{itemize}
  \item \textit{Typhoid Conjugate Vaccine (TCV):} E.g., Typbar-TCV. This is the preferred vaccine, consisting of the Vi capsular polysaccharide conjugated to a carrier protein (tetanus toxoid). It is administered as a single intramuscular injection. It can be given to infants as young as 6 months of age, provides long-lasting immunity (at least 5 years), and induces herd immunity.
  \item \textit{Vi Capsular Polysaccharide (ViPS) Vaccine:} An injectable subunit vaccine approved for individuals aged 2 years and older. It provides moderate protection (around 60\% to 70\%) and requires a booster dose every 2 to 3 years.
  \item \textit{Oral Ty21a Vaccine:} A live attenuated strain of \textit{S.} Typhi administered orally as 3 or 4 capsules on alternate days. It is approved for individuals aged 5 to 6 years and older and requires a booster every 5 years. It is contraindicated in pregnant women, immunocompromised individuals, and patients taking antibiotics.
  \item Note: Currently available vaccines target \textit{S.} Typhi only and do not provide reliable protection against \textit{S.} Paratyphi infection.
  \end{itemize}
\item \textbf{Management of Chronic Carriers:} Eradication of the carrier state is vital to prevent transmission. Carriers must be treated with a high-dose, prolonged course of antibiotics (such as ampicillin or ciprofloxacin for 4 to 6 weeks). If gallstones are present, surgical removal of the gallbladder (cholecystectomy) is often required in combination with antibiotics to successfully clear the infection. Carriers must be restricted from working in the food service industry or caring for children and elderly individuals until they have cleared the bacteria.
\item \textbf{Self-Care and Infection Control during Recovery:}
  \begin{itemize}
  \item Take all prescribed antibiotics exactly as directed by the physician, even if symptoms resolve earlier.
  \item Practice meticulous hand hygiene to prevent spreading the infection to family members and close contacts.
  \item Avoid preparing food for others until a doctor confirms that you are no longer shedding the bacteria. Clearance typically requires three consecutive negative stool cultures collected at least 24 hours apart, starting at least 48 hours after completing antibiotic therapy.
  \item Maintain adequate rest and nutrition, gradually transitioning from soft foods to a normal diet.
  \end{itemize}
\end{itemize}

\section*{Sources and Further Reading}
For authoritative information, clinical guidelines, and updates on enteric fever, refer to the following resources:
\begin{itemize}
\item \textbf{World Health Organization (WHO):} Typhoid fever health topic and fact sheet. Detailed global surveillance data, vaccine recommendations, and prevention strategies. \\ URL: \texttt{https://www.who.int/news-room/fact-sheets/detail/typhoid-fever}
\item \textbf{Centers for Disease Control and Prevention (CDC):} Typhoid and Paratyphoid Fever. Travel health notices, clinical description, and the CDC Yellow Book chapter for international travellers. \\ URL: \texttt{https://www.cdc.gov/typhoid-fever/index.html}
\item \textbf{The Lancet Infectious Diseases:} ``Enteric fever'' review articles. Comprehensive clinical, epidemiological, and microbiological overviews published by leading experts in international health. \\ URL: \texttt{https://www.thelancet.com/journals/laninf/home}
\item \textbf{UpToDate:} ``Epidemiology, microbiology, clinical manifestations, and diagnosis of enteric (typhoid and paratyphoid) fever.'' A peer-reviewed clinical decision support resource updated regularly by infectious disease specialists. \\ URL: \texttt{https://www.uptodate.com}
\item \textbf{Mayo Clinic:} Typhoid Fever. Consumer-friendly overview of symptoms, causes, diagnosis, treatment, and prevention. \\ URL: \texttt{https://www.mayoclinic.org/diseases-conditions/typhoid-fever/symptoms-causes/syc-20378661}
\item \textbf{National Health Service (NHS):} Typhoid Fever. Guidance on symptoms, treatment, vaccination, and travel advice from the United Kingdom's public health service. \\ URL: \texttt{https://www.nhs.uk/conditions/typhoid-fever/}
\item \textbf{Gavi, the Vaccine Alliance:} Typhoid Vaccine Support. Details on the deployment of typhoid conjugate vaccines (TCVs) in low-income, endemic countries. \\ URL: \texttt{https://www.gavi.org/types-support/vaccine-support/typhoid}
\end{itemize}